\chapter{Lyme disease}

\section*{Definition and Overview}
Laparoscopic or open exposure to tick-prone areas can lead to Lyme disease, which is a vector-borne infectious disease caused by the spirochete bacterium \textit{Borrelia burgdorferi} (and occasionally \textit{Borrelia mayonii}, \textit{Borrelia afzelii}, or \textit{Borrelia garinii}), transmitted to humans through the bite of infected blacklegged ticks (\textit{Ixodes} species). The disease is multisystemic, characterised by distinct stages: an early localised phase featuring a characteristic "bull's-eye" skin rash (erythema migrans), followed by early disseminated and late disseminated phases that can affect the joints, nervous system, and cardiovascular system. Lyme disease is endemic in parts of North America, Europe, and Asia. In Australia, local transmission remains a subject of ongoing clinical study, and the Australian Government recommends evaluating individuals who have travelled to endemic areas overseas.

\section*{Epidemiology}
Lyme disease is the most common tick-borne illness in the Northern Hemisphere. In the United States, approximately 30,000 cases are reported annually to the CDC, although estimates suggest the actual incidence is closer to 476,000 cases per year. In Europe, the estimated annual incidence is over 130,000 cases, with higher rates in Central and Northern European nations. Risk is seasonal, peaking during late spring and summer when nymphal ticks are active. In Australia, there is no conclusive scientific evidence of locally acquired classical Lyme disease due to the absence of confirmed vectors (\textit{Ixodes scapularis} or \textit{Ixodes pacificus}), but overseas travel to endemic regions accounts for all confirmed cases diagnosed within the country.

\section*{Aetiology and Pathophysiology}
Lyme disease is caused by infection with \textit{Borrelia} spirochetes transmitted during tick feeding. The pathophysiological mechanisms include:
\begin{itemize}
    \item \textbf{Bacterial transmission:} The spirochetes reside in the midgut of \textit{Ixodes} ticks and migrate to the salivary glands during feeding, a process requiring at least 24 to 36 hours of attachment.
    \item \textbf{Early localised hematogenous spread:} Outer surface proteins and bacterial lipoproteins interact with host tissue, facilitating local cutaneous migration (erythema migrans) and eventual vascular dissemination.
    \item \textbf{Immune evasion:} \textit{Borrelia} bacteria change their outer surface proteins (e.g. from OspA to OspC) and inhibit complement pathways, allowing them to persist in tissue niches.
    \item \textbf{Inflammatory response:} Infection triggers a host inflammatory response in target tissues (joints, meninges, myocardium), mediated by cytokines, leading to arthritis, neuropathies, and carditis.
    \item \textbf{Chronic sequelae:} Autoimmune cross-reactivity or persistent antigen debris may contribute to post-treatment symptoms in some individuals.
\end{itemize}

\section*{Clinical Features}
Lyme disease signs and symptoms depend on the stage of infection:
\begin{itemize}
    \item \textbf{Early localised stage (1--30 days):} Erythema migrans (a red, expanding rash often with central clearing, forming a "bull's-eye"), accompanied by flu-like symptoms such as fatigue, fever, headache, mild stiff neck, arthralgia, and myalgia.
    \item \textbf{Early disseminated stage (weeks to months):} Multiple secondary erythema migrans lesions, neurological manifestations (facial nerve palsy or Bell's palsy, lymphocytic meningitis, radiculopathy), and cardiac abnormalities (Lyme carditis presenting with fluctuating AV block).
    \item \textbf{Late disseminated stage (months to years):} Large joint arthritis (most commonly the knee) presenting with recurrent swelling and pain, and distal sensory neuropathies or encephalopathy.
    \item \textbf{Post-treatment Lyme disease syndrome (PTLDS):} Persistent fatigue, musculoskeletal pain, and cognitive difficulties continuing after standard antibiotic therapy.
\end{itemize}

\section*{Differential Diagnosis}
Because the symptoms are highly variable, clinicians must distinguish Lyme disease from other conditions:
\begin{itemize}
    \item \textbf{Southern tick-associated rash illness (STARI):} A tick-borne condition presenting with a similar rash but not caused by \textit{Borrelia burgdorferi}.
    \item \textbf{Rheumatoid arthritis:} Chronic autoimmune joint disease presenting with symmetric small joint involvement, unlike the asymmetric large joint involvement of Lyme arthritis.
    \item \textbf{Fibromyalgia:} Causes chronic widespread pain and fatigue without objective inflammatory signs or microbiological markers.
    \item \textbf{Multiple sclerosis:} Neurological symptoms like paresthesia and weakness can mimic Lyme neuroborreliosis.
    \item \textbf{Viral meningitis or Bell's palsy:} Idiopathic facial palsy or viral infections presenting similarly to early disseminated Lyme disease.
\end{itemize}

\section*{When to Seek Medical Care}
Individuals who have visited tick-prone areas (particularly overseas) and develop a rash, fever, or joint pain should consult a doctor.

\textbf{Call 000 (or local emergency services) for:}
\begin{itemize}
    \item \textbf{Cardiovascular instability:} Fainting (syncope), severe dizziness, palpitations, or chest pain, which can indicate high-grade AV block (Lyme carditis).
    \item \textbf{Severe neurological changes:} Sudden confusion, loss of consciousness, severe stiff neck with photophobia, or difficulty walking.
    \item \textbf{Difficulty breathing:} Shortness of breath or respiratory distress.
\end{itemize}
Other reasons to see a doctor include finding a tick attached to the skin that cannot be easily removed, or developing a spreading red rash within 3 to 30 days of a tick bite.

\section*{Diagnosis and Investigations}
Diagnosis is based on clinical signs, history of tick exposure in endemic regions, and validated laboratory tests:
\begin{itemize}
    \item \textbf{Clinical diagnosis:} A confirmed erythema migrans rash in a patient with a history of travel to an endemic area is diagnostic on its own, and antibiotic therapy should be started immediately without waiting for lab tests.
    \item \textbf{Two-tier serological testing:} The standard laboratory method. It consists of an initial enzyme-linked immunosorbent assay (ELISA), which, if positive or equivocal, is followed by a confirmatory Western blot test.
    \item \textbf{Polymerase chain reaction (PCR):} Used to detect \textit{Borrelia} DNA in joint fluid (synovial fluid) in cases of Lyme arthritis, or in skin biopsy specimens.
    \item \textbf{Lumbar puncture and CSF analysis:} Indicated in patients with suspected Lyme neuroborreliosis to check for intrathecal antibody production.
    \item \textbf{Electrocardiogram (ECG):} Performed to assess for conduction blocks in patients with suspected Lyme carditis.
\end{itemize}

\section*{Management}
Lyme disease is highly treatable with antibiotics, especially when diagnosed early:
\begin{itemize}
    \item \textbf{Pharmacological treatment (early stage):} Oral doxycycline (usually 100 mg twice daily for 10 to 14 days) is the first-line therapy. Amoxicillin or cefuroxime axetil are effective alternatives (especially for children or pregnant individuals).
    \item \textbf{Pharmacological treatment (disseminated stage):} Intravenous ceftriaxone (2 g daily for 14 to 28 days) is indicated for severe neurological manifestations or advanced Lyme carditis.
    \item \textbf{Symptomatic management:} NSAIDs are used to manage pain and inflammation in Lyme arthritis.
    \item \textbf{Tick removal technique:} Remove ticks immediately using fine-tipped tweezers by grasping the tick close to the skin and pulling straight up without twisting. Avoid squeezing the tick's body.
\end{itemize}

\section*{Complications}
Potential long-term complications of untreated or late-diagnosed Lyme disease include:
\begin{itemize}
    \item \textbf{Chronic Lyme arthritis:} Persistent joint inflammation and cartilage damage, primarily affecting the knees.
    \item \textbf{Neurological deficits:} Permanent facial weakness, distal sensory neuropathy, or mild cognitive impairment.
    \item \textbf{Permanent heart block:} Complete heart block requiring pacemaker implantation, though rare when treated.
    \item \textbf{Post-treatment Lyme disease syndrome:} Chronic fatigue, widespread pain, and brain fog that persist for months despite antibiotic therapy.
\end{itemize}

\section*{Prognosis and Outcomes}
The prognosis for early-stage Lyme disease treated with appropriate antibiotics is excellent, with almost all patients achieving complete recovery. For patients diagnosed in the late disseminated stage, response to therapy may be slower, but the infection is still successfully cleared in the vast majority of cases. A small subset of patients (around 5\% to 10\%) go on to develop Post-treatment Lyme disease syndrome; while there is no evidence that active bacterial infection persists in these cases, supportive care, physical therapy, and symptom management can help improve function over time.

\section*{Prevention and Self-Care}
Preventing tick bites is the primary way to prevent Lyme disease when visiting endemic areas:
\begin{itemize}
    \item \textbf{Wear protective clothing:} Wear long-sleeved shirts and long trousers tucked into socks when walking in wooded or grassy areas.
    \item \textbf{Use insect repellents:} Apply repellents containing DEET, picaridin, or IR3535 to exposed skin, and treat outdoor clothing with permethrin.
    \item \textbf{Perform tick checks:} Check the entire body, especially underarms, groin, and hair, for ticks immediately after returning from outdoor activities.
    \item \textbf{Shower promptly:} Shower within two hours of coming indoors to wash off unattached ticks.
    \item \textbf{Tick control in gardens:} Keep lawns mowed, clear brush, and create woodchip barriers between lawns and wooded areas to reduce tick habitats.
\end{itemize}

\section*{Sources and Further Reading}
The following reputable organisations provide further information on Lyme disease:
\begin{itemize}
    \item Australian Government Department of Health and Aged Care. \textit{Debilitating Symptom Complexes Attributed to Ticks (DSCATT) and Lyme disease.} https://www.health.gov.au/
    \item Centers for Disease Control and Prevention (CDC). \textit{Lyme disease symptoms, diagnosis, and treatment.} https://www.cdc.gov/
    \item World Health Organization (WHO). \textit{Vector-borne diseases: Lyme borreliosis.} https://www.who.int/
    \item Mayo Clinic. \textit{Lyme disease: patient education and clinical overview.} https://www.mayoclinic.org/
    \item NHS. \textit{Lyme disease information, tick bites, and prevention.} https://www.nhs.uk/
\end{itemize}
